% Sections 7.1 to 7.6, from lectures/L07/appendix/a_theory.md.

\section{Koordinatsystem och sträckor}
\label{c7:sec:koordinatsystem}
Ett \textbf{koordinatsystem} ($xy$-plan) har en horisontell $x$-axel och en vertikal $y$-axel
som möts i origo $(0, 0)$.

Sträckan $d$ mellan två punkter $(x_1, y_1)$ och $(x_2, y_2)$ ges av Pythagoras sats:
\[
  d = \sqrt{(x_2 - x_1)^2 + (y_2 - y_1)^2}
\]
Vinkeln som sträckan bildar med den positiva $x$-axeln är
\[
  \theta = \arctan\!\left(\frac{y_2 - y_1}{x_2 - x_1}\right).
\]

\begin{aside}
\note[Obs] Kontrollera alltid vilken kvadrant punkten befinner sig i och korrigera vinkeln vid
behov: $+180^{\circ}$ för kvadrant II och III.
\end{aside}

\section{Vektorer -- begrepp och notation}
\label{c7:sec:begrepp}
En \textbf{vektor} representerar en storhet med både storlek och riktning. I ett plan skrivs
vektorn $\mathbf{u}$ med sina komponenter:
\[
  \mathbf{u} = (u_x;\, u_y)
\]
Till skillnad från ett skalärt tal bär vektorn information om riktning. Vektorer används inom
elektroteknik för att representera spänningar och strömmar i AC-system.

\section{Absolutbelopp och vinkel}
\label{c7:sec:absolutbelopp}
\textbf{Absolutbeloppet} (längden) av $\mathbf{u} = (u_x;\, u_y)$ är
\[
  \abs{\mathbf{u}} = \sqrt{u_x^2 + u_y^2}.
\]
\textbf{Vinkeln} $v_u$ mot den positiva $x$-axeln är
\[
  v_u = \arctan\!\left(\frac{u_y}{u_x}\right) + \text{kvadrantkorrigering},
\]
där kvadrantkorrigeringen ges av tabellen.

\begin{narrowtable}{@{}cccl@{}}
  \thead{Kvadrant} & $u_x$ & $u_y$ & \thead{Korrigering} \\
  \midrule
  I & $+$ & $+$ & ingen \\
  II & $-$ & $+$ & $+180^{\circ}$ \\
  III & $-$ & $-$ & $+180^{\circ}$ \\
  IV & $+$ & $-$ & ingen (negativ vinkel) \\
\end{narrowtable}

\section{Vinkeln mellan två vektorer}
\label{c7:sec:vinkel}
Vinkeln $\alpha$ mellan $\mathbf{u}$ och $\mathbf{v}$ ges av
\[
  \cos \alpha = \frac{\mathbf{u} \cdot \mathbf{v}}{\abs{\mathbf{u}}\abs{\mathbf{v}}},
\]
där \textbf{skalärprodukten} är
\[
  \mathbf{u} \cdot \mathbf{v} = u_x v_x + u_y v_y.
\]

\section{Räkna med vektorer}
\label{c7:sec:rakna}

\subsection{Addition och subtraktion}
\label{c7:sec:addition}
\[
  \mathbf{u} \pm \mathbf{v} = (u_x \pm v_x;\; u_y \pm v_y)
\]

\subsection{Skalärmultiplikation}
\label{c7:sec:skalarmultiplikation}
\[
  k\mathbf{u} = (k u_x;\; k u_y)
\]
En negativ skalär byter riktning på vektorn. En motriktad enhetsvektor ges av
$-\dfrac{\mathbf{u}}{\abs{\mathbf{u}}}$.

\section{Typexempel}
\label{c7:sec:typexempel}

\begin{exempel}[Absolutbelopp och vinkel]
Vektorerna $\mathbf{u} = (3;\, 3)$ och $\mathbf{v} = (-2;\, 3)$ är givna.

\textbf{a)} Beräkna $\abs{\mathbf{u}}$ och $\abs{\mathbf{v}}$.
\begin{gather*}
  \abs{\mathbf{u}} = \sqrt{3^2 + 3^2} = \sqrt{18} = 3\sqrt{2} \approx 4{,}24 \\
  \abs{\mathbf{v}} = \sqrt{(-2)^2 + 3^2} = \sqrt{13} \approx 3{,}61
\end{gather*}

\textbf{b)} Beräkna vinklarna $v_u$ och $v_v$. Vektorn $\mathbf{u}$ ligger i kvadrant I:
\[
  v_u = \arctan\!\left(\frac{3}{3}\right) = 45^{\circ}
\]
Vektorn $\mathbf{v}$ ligger i kvadrant II, så $180^{\circ}$ läggs till:
\[
  v_v = \arctan\!\left(\frac{3}{-2}\right) + 180^{\circ} \approx -56{,}3^{\circ} + 180^{\circ}
  = 123{,}7^{\circ}
\]

\textbf{c)} Beräkna vinkeln mellan $\mathbf{u}$ och $\mathbf{v}$.
\begin{gather*}
  \mathbf{u} \cdot \mathbf{v} = 3 \cdot (-2) + 3 \cdot 3 = 3 \\
  \cos \alpha = \frac{3}{3\sqrt{2} \cdot \sqrt{13}} = \frac{3}{\sqrt{234}} \approx 0{,}196
  \quad \Rightarrow \quad \alpha \approx 78{,}7^{\circ}
\end{gather*}
\end{exempel}

\begin{exempel}[Vektorberäkning]
Beräkna $\mathbf{w} = \mathbf{u} + 2\mathbf{v}$ och $\abs{\mathbf{w}}$ för
$\mathbf{u} = (3;\, 3)$ och $\mathbf{v} = (-2;\, 3)$.
\begin{gather*}
  \mathbf{w} = (3;\, 3) + 2(-2;\, 3) = (3 - 4;\; 3 + 6) = (-1;\, 9) \\
  \abs{\mathbf{w}} = \sqrt{(-1)^2 + 9^2} = \sqrt{82} \approx 9{,}06
\end{gather*}
\end{exempel}

% Keeps this example from starting at the foot of a page, where it is drawn over the running
% head of the next.
\needspace{8\baselineskip}
\begin{exempel}[Motriktad vektor med given längd]
Bestäm en vektor med längden $5$ som är motriktad $\mathbf{u} = (3;\, 3)$.
\[
  \mathbf{w} = -\frac{5}{\abs{\mathbf{u}}}\,\mathbf{u} = -\frac{5}{3\sqrt{2}}(3;\, 3)
  = \left(-\frac{5}{\sqrt{2}};\; -\frac{5}{\sqrt{2}}\right) \approx (-3{,}54;\; -3{,}54)
\]
\end{exempel}
